% !PRINTED 524
\weberpage{524}
%<<CONT
and
\weq{7}{\begin{gathered}
v = \alpha + \varepsilon\alpha_1 + \varepsilon^2\alpha_2 \\
v' = \alpha + \varepsilon^2\alpha_1 + \varepsilon\alpha_2;
\end{gathered}}
then the effect of $\pi$ is that $v$ passes into $\varepsilon^2 v$, and $v'$ into $\varepsilon v'$, while $v^3$, $v'^3$ and $vv'$ remain unchanged, and hence are contained in the field $\Omega'$. There is no difficulty in computing these quantities. If for brevity we set
\[
\begin{gathered}
A = \alpha^2\alpha_1 + \alpha_1^2\alpha_2 + \alpha_2^2\alpha, \\
A' = \alpha_1^2\alpha + \alpha_2^2\alpha_1 + \alpha^2\alpha_2,
\end{gathered}
\]
then
\[
\begin{gathered}
A + A' = ab - 3c \\
A - A' = \sqrt{D},
\end{gathered}
\]
and for $v^3$ one obtains
\[
v^3 = \alpha^3 + \alpha_1^3 + \alpha_2^3 + 6\alpha\alpha_1\alpha_2 + 3\varepsilon A + 3\varepsilon^2 A',
\]
hence
\weq{8}{v^3 = a^3 - \frac{9}{2}ab + \frac{27}{2}c + \frac{3}{2}\sqrt{-3D},}
and likewise
\weq{9}{v'^3 = a^3 - \frac{9}{2}ab + \frac{27}{2}c - \frac{3}{2}\sqrt{-3D}.}

But we also have
\weq{10}{vv' = \alpha^2 + \alpha_1^2 + \alpha_2^2 - \alpha\alpha_1 - \alpha\alpha_2 - \alpha_1\alpha_2 = a^2 - 3b,}
and from (8), (9), (10) one obtains, in agreement with an earlier formula [\S. 47, (8)],
\[
27D = 4(a^2 - 3b)^3 - (2a^3 - 9ab + 27c)^2.
\]
If to (7) we further add the equation
\[
a = \alpha + \alpha_1 + \alpha_2,
\]
then one finds, in agreement with the \gesperrt{Cardan} formula,
\[
\begin{gathered}
3\alpha = a + v + v' \\
3\alpha_1 = a + \varepsilon^2 v + \varepsilon v' \\
3\alpha_2 = a + \varepsilon v + \varepsilon^2 v'.
\end{gathered}
\]

\begin{center}
§. 160.

\gesperrt{Permutation groups of four elements.}
\end{center}

A good example of the \gesperrt{Galois} theory is offered by the equations of the fourth degree, where the circumstances are so simple
%CONT>>
